% !TeX root = ../papers/nc-p3.tex
\section{被約窓での階数一曲面の排除と階数二平面上の曲線}
\label{sec:rank-one-surface-exclusion}
\begin{theorem}\label{thm:rank-one-surface-exclusion}
四辺有限の標準有限Frobenius窓で $X$ は被約、全二次式を成分に持つ反対平方根 $N$ は一般可逆とする。
階数一軌跡は高々曲線である。
\end{theorem}
\begin{proof}[階数一曲面からの一次表示]
階数一の既約曲面 $S$ があれば、その一般点で単元一行を消すことにより $S$ の三重が八次行列式に入る。
従って $S$ は平面か既約二次曲面である。$N|_S$ はVeroneseを通じて射 $j:S\to\mathbb P^3$ を与える。
$j^*\mathcal O(2)=\mathcal O_S(2)$ で、これらの曲面のPicard群に二捩れはない
（錐は定理\ref{thm:low-rank-surface-factors}の解消を使う）。従って $j$ は周囲の一次式の列 $a$ で書け、
定数を吸収して $N|_S=aa^t$ となる。$j$ は閉埋込みである。実際 $N$ の射は完全二次Veroneseの基底変更である。
$S$ が二次曲面なら $N=aa^t+qC$ で $C$ は定数対称行列である。二次評価が単射なので
$q\ne0$ となる印が存在し、その印で二つの階数一以下の行列の差を取ると $\operatorname{rank}C\le2$。
このとき $N$ は一般点でも階数3以下となり矛盾する。

$S$ が平面 $h=0$ なら $N=aa^t+hC$ で $C$ は一次対称行列である。五行の一次対称行列と五次式
\[
              A=\begin{pmatrix}C&a\\a^t&-h\end{pmatrix},
                 \qquad F=\det A=-\det N/h^3
\]
を取る。$h\ne0$ では $\operatorname{rank}A=1+\operatorname{rank}N$。
従って平面外の $n\ge13$ 個の印全てで $A$ の階数は2であり $F$ は三重に消える。
\end{proof}
\begin{proof}[残余五次式の既約性]
$F$ の各因子について平面外の $n$ 印での重複度を加算する。
一次・二次因子の上界7と15よりこれらだけの分割では総和は高々 $15+15+7=37<3n$。
既約三次の寄与は前定理の特異点の評価で $n+4$ 以下、残る次数二は15以下でやはり $3n$ に届かない。
既約四次と一次の分割も不可能である。四次の二重以上は高々 $n$ 印、三重以上は共線で三印以下、
四重は高々一印である。従って増分を含む四次の寄与は $2n+4$ 以下、一次を足しても $2n+11<3n$。
これで $F$ は既約かつ被約である。
\end{proof}
\begin{proof}[二つの階数二点の割線]
階数二以下の二点を結ぶ入力直線では、$A$ が一次なので階数は高々4であり、直線は $V(F)$ に入る。
従って階数二以下の軌跡は有限である。実際、非平面曲線なら割線が三空間を覆い、非直線の平面曲線なら
割線が平面を覆う。直線成分があってもその上の印は三つ以下で、別の階数二の印とその直線を結ぶ平面が
$V(F)$ に入る。いずれも既約五次に反する。また階数三以下の軌跡は高々曲線である。
曲面成分ならその二重が $F$ に入ってしまうからである。
一般の三次元定数部分空間への制限 $3\mathcal O(-1)\to5\mathcal O$ の階数二以下の軌跡は有限となる。
実際階数四の層上ではその核直線を含む条件が余次元2、階数三の層上では核二平面と交わる条件が余次元1で、
他の退化は有限な階数二以下の層に限られる。期待余次元の退化長は
$[H^3](1-H)^{-3}=10$ だが、階数二の指定印だけで十三個以上あり矛盾する。
\end{proof}

\begin{proposition}\label{prop:rank-two-plane-conic-trace}
同じ被約窓で一般階数二の入力平面上に階数一の曲線があれば、その曲線は重複を許す平面二次式 $q$ の零集合である。
平面上の行列は、定数出力列 $v$ と一次列 $a$ について $aa^t+q vv^t$ と書ける。
$v$ による商への $a$ の射影は入力平面から商の射影平面への線型同型である。
\end{proposition}
\begin{proof}
平面の三変数環で第二外冪を $d gg^t$ と原始に書く。一般階数二より $g$ は像二平面のPlücker列である。
$\deg g\le2$、次数2なら $d$ は定数で原始列の基点は有限となり階数一曲線はない。
定数なら像は固定で全成分の包は高々三次元となり、平面上の全二次式が六次元であることに反する。
従って $g$ は一次式の列で、その線型像はKlein二次曲面内に入る。
像が直線以下なら二平面の和は固定三空間以下であり、対応する固定列 $Nv$ は入力平面の式で割れるので
定理\ref{thm:opposite-column-divisor}に反する。像平面は、全ての二平面が共通の直線を含むものか、
固定三空間に入るものかのいずれかである。実際Plücker直線で結ばれる射影直線は交わり、
二つを固定すれば他のものはその交点を通るかその張る平面に入る。線型像の既約性より一方が全体で成立する。
後者は同じ固定列で除かれる。前者で共通ベクトルを第一座標 $v$ にし、像を $\langle v,(0,x)^t\rangle$
と書く。$x$ の三成分は平面の独立な座標である。下三行三列は $cxx^t$、上行の混合部は $l x^t$ である。
これは一般像への包含と $x$ の原始性による。次数より $c$ は定数、$l$ は一次式である。
$c=0$ なら非零成分は高々四つで不可能であり、平方完成で主張を得る。
二次式 $q$ は一般階数二から非零であり、$x$ の一成分が単元の各開集合で階数一イデアルは $(q)$ である。
\end{proof}
